\scanpage{73}
% archive sheet number: 72

\[
\underline{G}(k')\handwritten{\longrightarrow}\underline{G}(k)\ .
\]
Ce dernier rend commutatif \inserted{(mod isomorphisme canonique)} le diagramme
\diagram{\begin{tikzcd}[ampersand replacement=\&,column sep=large,row sep=large]
\underline{G}(k\handwritten{'}) \handwritten{\arrow[r]} \handwritten{\arrow[d,dash,"{\unsure{\ }}" description]} \& \underline{G}(k^{\deleted{\unsure{'}}}) \handwritten{\arrow[d,dash,"{\unsure{\ }}" description]} \\
\handwritten{\unsure{\underline{\pi}|\underline{\Pi}}_{p}}(k')\mathrm{x}\handwritten{[}G_m\handwritten{]}_{\underline{\unsure{\mathbb{Q}|Q}}_p} \handwritten{\arrow[r]} \& \handwritten{\underline{\pi}}_p(k)\mathrm{x}\handwritten{[}\underline{G}_m\handwritten{]}_{\underline{Q}_p}
\end{tikzcd}}
(i.e. on a, sur les catégories tannakiennes\deleted{\unsure{, un}} correspondantes, un dia\unsure{g|gg}ramme comm\unsure{u|y}tatif à isomorphisme \inserted{can.} près).

Oubliant dans ce dernier les facteurs $\handwritten{[}\underline{G}_m\handwritten{]}$, on voit donc que $\underline{G}(k')\handwritten{\longrightarrow}\underline{G}(k)$ se factorise par la sous-\deleted{catégorie} gerbe de $\underline{G}(k)$ correspondant au sous-groupe fermé de $\handwritten{\pi}_1(k)$ image de $\handwritten{\pi}_1(k')$ ... A prouver:

\uline{Proposition} \handwritten{6.1.} Si $k'$ est une extension \deleted{finie} étale de $k$, alors l'homomorphisme de gerbes $\underline{G}(k')\handwritten{\longrightarrow}\underline{G}(k)$ est un \uline{monomorphisme} dont l'image est exactement la sous-gerbe ``ouverte'' définie ci-dessus. En particulier, si $k'$ est une extension finie galoisienne de $k$, \deleted{on a une} de groupe de Galois $g$, et on une suite exacte canonique de gerbes
\[
0\handwritten{\longrightarrow}\underset{\handwritten{\deleted{\unsure{\pi|\mu}}}}{\underline{G}^{\deleted{\unsure{'}}}}(k\handwritten{'})\handwritten{\longrightarrow}\underline{G}(k)\handwritten{\longrightarrow [}g\handwritten{]}_{\underline{Q}_p}\handwritten{\longrightarrow}0\ .
\]
Cet énoncé devrait être formellement équivalent au suivant: $\operatorname{Spec}(k')\handwritten{\longrightarrow}\operatorname{Spec}(k)$ est un morphisme de descente effective pour la catégorie fibrée des $\mathrm{Fcriso}(L)$ sur des composés finis\inserted{$L$} de corps de car. $p$. La démonstration de ce dernier fait est $\handwritten{\pm}$ triviale.

\uline{Question} \handwritten{6.2.} A-t-on encore un énoncé analogue lorsqu'on suppose que $k'$ est une extension galoisienne de $k$ qui peut-être infinie ? Le cas le plus intéressant serait celui où $k'$est la clôture séparable de $k$, de sorte qu'on aurait un dévissage de $\underline{G}(k)$ en
\[
0\handwritten{\longrightarrow}\underline{G}(\handwritten{\bar{k}}^{\deleted{\unsure{'}}})\handwritten{\longrightarrow}\underline{G}(k)\handwritten{\longrightarrow\underline{\pi}}_p(k)\handwritten{\longrightarrow}0\quad\deleted{\unsure{?}}\ ;
\tag*{\handwritten{(\unsure{6}.2.1)}}
\]
lorsque $k$ est parfait i.e. $\handwritten{\bar{k}}$ alg. clos, nous verrons que $\underline{G}(\handwritten{\bar{k}})$ ne dépend pas essentiellement de \handwritten{$k$. À}
